\documentclass[10pt,a4paper]{amsart}
\usepackage[margin=19mm]{geometry}
\usepackage{amsmath,amssymb,mathtools}
\usepackage[scr=rsfs,cal=euler]{mathalfa}
\usepackage{xfrac}
\usepackage[unicode,hidelinks,bookmarks=false]{hyperref}
\newcommand{\ie}{i.e.\ }
\newcommand{\cf}{cf.\ }
\newcommand{\resp}{respectively\ }
\newcommand{\Subsection}{\subsection}
\providecommand{\nobreakdash}{\nobreak\hbox{-}}
\setlength{\emergencystretch}{2em}
\multlinegap0pt
\usepackage{bm}
\usepackage{bbm}
\usepackage{dsfont}
\usepackage{xspace}
\usepackage{cancel}
\usepackage{mathtools}
\usepackage{cleveref}
\usepackage{amsthm}
\makeatletter
\@ifundefined{lemma}{\newtheorem{lemma}{Lemma}}{}
\makeatother
\title[Minimax Private Estimation of Smooth Optimal-Transport Maps]{Minimax Private Estimation of Smooth Optimal-Transport Maps}
\author{Cl\'ement Lalanne, David Rodr\'iguez-V\'itores, Franck Iutzeler, Jean-Michel Loubes}
\date{Source: arXiv:2606.04683v1 (CC BY 4.0)}
\begin{document}
\maketitle

\noindent\emph{Source and licence.} Minimax Private Estimation of Smooth Optimal-Transport Maps. Version arXiv:2606.04683v1 (CC BY 4.0), distributed under the Creative Commons Attribution 4.0 International licence, as recorded in the arXiv licence metadata for this exact version and shown on the abstract page https://arxiv.org/abs/2606.04683v1. Full rights notice: licenses/arxiv-2606.04683-CC-BY-4.0.txt.\par\smallskip

\noindent\textit{Editorial scope.} Counted unit: the Lemma whose source label is \texttt{lemma:holder\_inversion} in the article (source0.tex lines 2250--2261, source label \texttt{lemma:holder\_inversion}), delivered together with the article's own complete original proof of it (source0.tex lines 2262--2457, one \texttt{proof} environment of 196 lines). No mathematics is written, repaired or re-derived: statement and proof are the article's own text line for line. Required context retained uncounted: lemma:holder\_composition (lines 2176--2182). Every definition, display and label of those lines is retained unchanged. Omitted from this package and not redistributed: 1--2175, 2183--2249, 2458--2833. Nothing inside a retained excerpt is omitted or altered: every retained body is delivered exactly as the article's lines are written, so no omission receipt of any kind accompanies this package. Citations: the retained excerpts carry no citation command at all, so no bibliography entry of the article is transcribed and no reference is redistributed. Reference adaptation: none was needed and none was made; every reference inside the delivered unit resolves inside it, so no unresolved reference or citation is produced. Printed slips kept literal: no printed word, symbol, hypothesis or number of the counted statement or of its proof was altered. Layout and class: the article's own class is \texttt{article} with packages including placeins, color-edits, xcolor, enumitem, neurips\_2026, inputenc, fontenc, hyperref, url, booktabs, amsfonts, nicefrac, microtype, amsmath; the wrapper is the lane's pinned \texttt{amsart} editorial wrapper at 10pt on A4, which restates the theorem-like environments the retained text needs and adds the article's own macro definitions that the retained text needs (none), with extra wrapper packages (bm, bbm, dsfont, xspace, cancel, mathtools, cleveref). The delivered numbering is the unit's own. Assets: no retained span contains a figure environment, a graphics-inclusion command, a PDF-page inclusion, a table or a TikZ picture; no image file is copied into this package, the package declares no asset dependency and no image file is redistributed with it. Licence: version arXiv:2606.04683v1 of this article is distributed under the Creative Commons Attribution 4.0 International licence, as recorded in the arXiv licence metadata for this exact version and shown on the abstract page https://arxiv.org/abs/2606.04683v1. A full rights notice is delivered at licenses/arxiv-2606.04683-CC-BY-4.0.txt. Characters: every retained excerpt is pure ASCII, so the delivered engine, XeLaTeX with its default Latin Modern text font, reports no missing character.
\begin{lemma}
\label{lemma:holder_composition}
    If $u$ and $v$ are in $C^{\infty}(\Omega)$ with $v(\Omega) \subseteq \Omega$ and $s \geq 0$, then 
    \begin{equation}
        \| u \circ v\|_{C^{s}(\Omega)} \lesssim \| u \|_{C^{s}(\Omega)} (1 + \max (\| v\|_{C^{s}(\Omega)}^s, \| v\|_{C^{1}(\Omega)}^s)) \;.
    \end{equation} 
\end{lemma}

\begin{lemma}
    \label{lemma:holder_inversion}
    Using the same notations as in the proof, if $\nabla \phi_\theta$ satisfies
    \begin{equation}
    \|\nabla \phi_\theta - \mathrm{Id}\|_{C^{s+1}(\Omega)} \leq \epsilon,
\end{equation}
for a small enough $\epsilon$ (depending on $s$, $d$ and $\Omega$), then 
\begin{equation}
    \|\nabla \phi_\theta^{-1}\|_{C^{s+1}(\Omega)} \leq C
\end{equation}
where $C$ is a constant depending only on $d$, $s$ and $\Omega$.
\end{lemma}

\begin{proof}
   Write
\begin{equation}
\nabla \phi_\theta(x) = x + w_\theta(x),
\qquad
w_\theta := \nabla \phi_\theta - \mathrm{Id}.
\end{equation}
By assumption,
\begin{equation}
\label{eq:ass_w_small}
\|w_\theta\|_{C^{s+1}(\Omega)} \le \epsilon.
\end{equation}

Since
$
D(\nabla\phi_\theta)(x) = I + Dw_\theta(x),
$ 
we have
$
\|Dw_\theta\|_{L^\infty(\Omega)} \le \|w_\theta\|_{C^{s+1}(\Omega)} \le \epsilon
$ by definition of the Holder norms.
Assume $\epsilon\le 1/2$. Then for all $x\in\Omega$,
\begin{equation}
\|D(\nabla\phi_\theta)(x) - I\| \le \frac12,
\end{equation}
hence $D(\nabla\phi_\theta)(x)$ is invertible and
\begin{equation}
\label{eq:jac_inv_bound}
\bigl\|D(\nabla\phi_\theta)(x)^{-1}\bigr\| \le \frac{1}{1-\|D(\nabla\phi_\theta)(x)-I\|} \le 2.
\end{equation}
In particular, $\nabla\phi_\theta$ is a $C^1$ diffeomorphism of $\Omega$ onto itself and, by the inverse function theorem,
$\nabla\phi_\theta^{-1}$ is $C^1$. Moreover, for all $y\in\Omega$,
\begin{equation}
\label{eq:Ds_formula}
D(\nabla\phi_\theta^{-1})(y)
=
\Bigl(D(\nabla\phi_\theta)\bigl(\nabla\phi_\theta^{-1}(y)\bigr)\Bigr)^{-1},
\end{equation}
so by \eqref{eq:jac_inv_bound},
\begin{equation}
\label{eq:lip_inverse}
\|D(\nabla\phi_\theta^{-1})\|_{L^\infty(\Omega)} \le 2,
\end{equation}
i.e.\ $\nabla\phi_\theta^{-1}$ is uniformly Lipschitz.

We will work from the implicit equation that defines the inverse.
For all $y\in\Omega$,
\begin{equation}
\label{eq:implicit_identity}
\nabla\phi_\theta\bigl(\nabla\phi_\theta^{-1}(y)\bigr) = y.
\end{equation}
Using $\nabla\phi_\theta = \mathrm{Id}+w_\theta$, this rewrites as
\begin{equation}
\label{eq:implicit_identity_2}
\nabla\phi_\theta^{-1}(y) + w_\theta\bigl(\nabla\phi_\theta^{-1}(y)\bigr) = y.
\end{equation}

We now look at integer derivatives up to order $\lfloor s+1\rfloor$.
Let
$
m := \lfloor s+1\rfloor.
$
Let $\alpha$ be a multi-index with $1\le |\alpha| \le m$. Differentiating \eqref{eq:implicit_identity_2} yields
\begin{equation}
\label{eq:diff_identity_alpha}
\partial^\alpha \nabla\phi_\theta^{-1}(y) + \partial^\alpha\!\Bigl(w_\theta\circ\nabla\phi_\theta^{-1}\Bigr)(y) = 0.
\end{equation}
By the Fa\`a di Bruno formula, $\partial^\alpha(w_\theta\circ\nabla\phi_\theta^{-1})$ is a finite sum of terms
\begin{equation}
\label{eq:fdb_general_term}
\partial^k w_\theta\bigl(\nabla\phi_\theta^{-1}(y)\bigr)\,
\prod_{j=1}^k \partial^{\gamma_j}\nabla\phi_\theta^{-1}(y),
\qquad
1\le k\le|\alpha|,
\qquad
\sum_{j=1}^k |\gamma_j| = |\alpha|.
\end{equation}
Among these terms, the ones with $k=1$ and $|\gamma_1|=|\alpha|$ equal
$Dw_\theta(\nabla\phi_\theta^{-1}(y))\,\partial^\alpha \nabla\phi_\theta^{-1}(y)$.
Thus \eqref{eq:diff_identity_alpha} can be rewritten as
\begin{equation}
\label{eq:solve_alpha}
\Bigl(I + Dw_\theta\bigl(\nabla\phi_\theta^{-1}(y)\bigr)\Bigr)\,\partial^\alpha \nabla\phi_\theta^{-1}(y)
=
- R_\alpha(y),
\end{equation}
where $R_\alpha(y)$ is a finite sum of terms of the form \eqref{eq:fdb_general_term} with $k\ge 2$, hence involving only
derivatives of $\nabla\phi_\theta^{-1}$ of order $<|\alpha|$.

Since $\|Dw_\theta\|_{L^\infty}\le\epsilon\le 1/2$, we have
\begin{equation}
\label{eq:inv_matrix_bound}
\sup_{y\in\Omega}\left\|\Bigl(I+Dw_\theta(\nabla\phi_\theta^{-1}(y))\Bigr)^{-1}\right\|
\le 2.
\end{equation}
Taking $L^\infty$ norms in \eqref{eq:solve_alpha} gives
\begin{equation}
\label{eq:alpha_Linf_bound}
\|\partial^\alpha \nabla\phi_\theta^{-1}\|_{L^\infty(\Omega)}
\le
2\,\|R_\alpha\|_{L^\infty(\Omega)}.
\end{equation}
We now argue by induction on $l:=|\alpha|$. The case $l=2$ follows from \eqref{eq:lip_inverse}.
Assume that for some $l\in\{2,\dots,m\}$,
\begin{equation}
\label{eq:induction_hyp}
\max_{|\gamma|\le l-1}\|\partial^\gamma \nabla\phi_\theta^{-1}\|_{L^\infty(\Omega)} \le C_{l-1},
\end{equation}
and let $|\alpha|=l$. We then obtain
\begin{equation}
\|R_\alpha\|_{L^\infty(\Omega)}
\lesssim
\sum_{k=2}^l
\|\partial^k w_\theta\|_{L^\infty(\Omega)}\,
C_{l-1}^{\,k}
\lesssim
\|w_\theta\|_{C^l(\Omega)}\,C_{l-1}^{\,l}
\le
\epsilon\,C_{l-1}^{\,l}.
\end{equation}
Plugging this into \eqref{eq:alpha_Linf_bound} yields
\begin{equation}
\max_{|\alpha|=l}\|\partial^\alpha \nabla\phi_\theta^{-1}\|_{L^\infty(\Omega)}
\le
C(d,l,\Omega)\,\epsilon\,C_{l-1}^{\,l}.
\end{equation}
Choosing $\epsilon>0$ small enough (depending only on $d,s,\Omega$) and iterating over $l=1,\dots,m$ yields a uniform bound
\begin{equation}
\label{eq:Cm_bound_inverse}
\max_{|\alpha|\le m}\|\partial^\alpha \nabla\phi_\theta^{-1}\|_{L^\infty(\Omega)} \le C.
\end{equation}

For the non-integer part, let
$
r := (s+1) - m \in [0,1).
$
If $r=0$, we are done. Otherwise fix a multi-index $\alpha$ with $|\alpha|=m$.
From \eqref{eq:solve_alpha} we have
\begin{equation}
\label{eq:alpha_product_clean}
\partial^\alpha \nabla\phi_\theta^{-1}(y)
=
-\Bigl(I + Dw_\theta(\nabla\phi_\theta^{-1}(y))\Bigr)^{-1} R_\alpha(y).
\end{equation}
Using the $C^r$ product rule (as in the proof of \Cref{lemma:holder_composition}), we obtain
\begin{equation}
\label{eq:Cr_product_bound}
\begin{aligned}
&\|\partial^\alpha \nabla\phi_\theta^{-1}\|_{C^r(\Omega)} \\
&\lesssim
\left\|\Bigl(I+Dw_\theta\circ\nabla\phi_\theta^{-1}\Bigr)^{-1}\right\|_{L^\infty(\Omega)}\|R_\alpha\|_{C^r(\Omega)}
+
\left\|\Bigl(I+Dw_\theta\circ\nabla\phi_\theta^{-1}\Bigr)^{-1}\right\|_{C^r(\Omega)}\|R_\alpha\|_{L^\infty(\Omega)}.
\end{aligned}
\end{equation}
The first factor is bounded by $2$ by \eqref{eq:inv_matrix_bound}. We now bound the remaining terms.

First, each term in $R_\alpha$ is a product of factors of the form $\partial^k w_\theta\circ\nabla\phi_\theta^{-1}$
(with $2\le k\le m$) and $\partial^\gamma \nabla\phi_\theta^{-1}$ with $|\gamma|\le m-1$.
Iterating the same $C^r$ product rule and using \Cref{lemma:holder_composition} together with
\eqref{eq:ass_w_small} and the uniform $C^m$ bound \eqref{eq:Cm_bound_inverse} yields
\begin{equation}
\label{eq:Ralpha_Cr_bound}
\|R_\alpha\|_{C^r(\Omega)} + \|R_\alpha\|_{L^\infty(\Omega)} \le C\,\epsilon.
\end{equation}

Second, since $A\mapsto(I+A)^{-1}$ is smooth on $\{A:\|A\|\le 1/2\}$, we have
\begin{equation}
\label{eq:inverse_map_bound}
\left\|\Bigl(I+Dw_\theta\circ\nabla\phi_\theta^{-1}\Bigr)^{-1}\right\|_{C^r(\Omega)}
\lesssim
1 + \|Dw_\theta\circ\nabla\phi_\theta^{-1}\|_{C^r(\Omega)}.
\end{equation}
Moreover, by \Cref{lemma:holder_composition},
\begin{equation}
\label{eq:comp_bound_Dw}
\|Dw_\theta\circ\nabla\phi_\theta^{-1}\|_{C^r(\Omega)}
\lesssim
\|Dw_\theta\|_{C^r(\Omega)}\Bigl(1+\|\nabla\phi_\theta^{-1}\|_{C^{1}(\Omega)}^{r}\Bigr)
\lesssim
\|w_\theta\|_{C^{s+1}(\Omega)}
\le \epsilon,
\end{equation}
where we used \eqref{eq:lip_inverse} to control $\|\nabla\phi_\theta^{-1}\|_{C^1}$.
Plugging \eqref{eq:Ralpha_Cr_bound}--\eqref{eq:comp_bound_Dw} into \eqref{eq:Cr_product_bound} gives
\begin{equation}
\max_{|\alpha|=m}\|\partial^\alpha \nabla\phi_\theta^{-1}\|_{C^r(\Omega)} \le C.
\end{equation}


Combining \eqref{eq:Cm_bound_inverse} with the previous result yields
\begin{equation}
\|\nabla\phi_\theta^{-1}\|_{C^{s+1}(\Omega)} \le C,
\end{equation}
where $C$ depends only on $d$, $s$ and $\Omega$.
\end{proof}

\end{document}
